\documentclass[11pt,letterpaper]{article}
\PassOptionsToPackage{hyphens}{url}
\usepackage[margin=.83in,top=.8in,bottom=.8in,headheight=15pt]{geometry}
\usepackage[T1]{fontenc}
\usepackage{lmodern,amsmath,amssymb,mathtools,bm,booktabs,array,graphicx,tikz}
\usetikzlibrary{arrows.meta,positioning,calc,decorations.pathreplacing,patterns}
\definecolor{ink}{HTML}{17344D}\definecolor{accent}{HTML}{087E8B}\definecolor{pale}{HTML}{EDF5F5}\definecolor{gold}{HTML}{AD6B19}
\providecommand{\dd}{\,\mathrm{d}}\providecommand{\R}{\mathbb{R}}
\providecommand{\E}{\mathbb{E}}\providecommand{\Var}{\operatorname{Var}}
\providecommand{\Cov}{\operatorname{Cov}}\providecommand{\diag}{\operatorname{diag}}
\providecommand{\norm}[1]{\left\lVert#1\right\rVert}
\providecommand{\abs}[1]{\left|#1\right|}
\usepackage{microtype,hyperref,fancyhdr,enumitem,amsthm,sectsty,needspace}
\allsectionsfont{\raggedright}
\hypersetup{colorlinks=true,linkcolor=ink,urlcolor=accent,pdfauthor={Zan Ahmad},pdftitle={Population transfers: SIR dynamics, thresholds, and numerical conservation / Review}}
\setlength{\parindent}{0pt}\setlength{\parskip}{6pt}\linespread{1.07}
\setlist{itemsep=3pt,topsep=4pt,leftmargin=*}
\emergencystretch=2em\clubpenalty=10000\widowpenalty=10000
\pagestyle{fancy}\fancyhf{}\fancyhead[L]{\small Mathematical Modeling}\fancyhead[R]{\small Lecture 6}
\fancyfoot[L]{\footnotesize Zan Ahmad}\fancyfoot[R]{\footnotesize\thepage}\renewcommand{\headrulewidth}{.3pt}
\newtheorem{proposition}{Proposition}\newtheorem{theorem}{Theorem}\newtheorem{lemma}{Lemma}
\theoremstyle{definition}\newtheorem{definition}{Definition}\newtheorem{example}{Example}
\numberwithin{equation}{section}
\begin{document}
\section*{Population transfers: SIR dynamics, thresholds, and numerical conservation --- review sheet}
Use this after attempting the worked examples and practice. Explain each claim in words before calculating.

\section*{Key ideas}
\begin{itemize}
\item Internal transfers change group composition while conserving a closed total.
\item Infection flow $\beta$SI/N and recovery flow $\gamma$I have units people/day.
\item The initial direction of I follows the comparison $\beta$S/N versus $\gamma$.
\item A sufficient nonnegative-step condition does not guarantee numerical accuracy.
\item Synthetic scenarios and validated forecasts require different evidence.
\end{itemize}
\section*{Common errors and repairs}
\textbf{Subtracting recovered people from total N.} R remains inside this closed population; recovery moves people between categories.

\textbf{Calling $\beta$ a probability with no time unit.} In this model $\beta$ is an effective rate in 1/day; a transfer also depends on current group sizes and interval length.

\textbf{Treating a constant total as proof that the entire trajectory is correct.} Check bounds, limiting cases and step refinement as well as conservation.

\section*{Retrieve without your notes}
\begin{itemize}
\item Which two changes cancel for an infection transfer?
\item Can a positive $\beta$ start an outbreak when I0 =\allowbreak{} 0 in this closed model?
\item Why compare solutions at common times when changing the step?
\end{itemize}
\section*{Readiness check}
\begin{itemize}
\item Draw a closed-population transfer diagram with clearly defined categories.
\item Compute infection and recovery transfers over a short interval.
\item Check conservation and nonnegative group sizes in a numerical update.
\item Predict whether the infectious group initially increases from the competing rates.
\item Distinguish a controlled synthetic scenario from a validated population forecast.
\end{itemize}
\end{document}
